% ============================================================================
%  Rozdział: Proces tworzenia syntetycznego datasetu kalibracyjnego
%  Plik przeznaczony do samodzielnej kompilacji lub włączenia przez \input{}
%  do głównego dokumentu pracy magisterskiej.
% ============================================================================
\documentclass[12pt,a4paper]{report}
\usepackage[utf8]{inputenc}
\usepackage[T1]{fontenc}
\usepackage[polish]{babel}
\usepackage{amsmath,amssymb}
\usepackage{booktabs}
\usepackage{array}
\usepackage{geometry}
\geometry{margin=2.5cm}

\begin{document}

\chapter{Proces tworzenia syntetycznego datasetu kalibracyjnego}
\label{chap:tworzenie-datasetu}

\section{Cel generowania danych syntetycznych}
Proces tworzenia datasetu polega na wygenerowaniu zbioru obrazów przedstawiających płaską
szachownicę obserwowaną przez syntetyczną kamerę o znanych parametrach wewnętrznych i znanych
współczynnikach dystorsji. Tak przygotowany zbiór danych ma zastąpić lub uzupełnić klasyczną
procedurę ręcznego zbierania obrazów kalibracyjnych. Najważniejsza różnica polega na tym, że
w zbiorze syntetycznym wartości parametrów kamery są znane dokładnie, ponieważ są one losowane
przed wygenerowaniem obrazu. Dzięki temu każdy obraz lub każda sekwencja obrazów może zostać
opisana etykietą regresyjną, którą model neuronowy ma później odtworzyć.

W tym rozdziale opisano proces generowania danych od strony matematycznej. Nie jest istotna
struktura programu ani sposób organizacji plików, lecz ciąg przekształceń geometrycznych,
który prowadzi od idealnego wzorca szachownicy do końcowego obrazu widzianego przez kamerę.
Cały proces można potraktować jako złożenie czterech głównych operacji:
\begin{enumerate}
    \item zdefiniowania płaskiego wzorca w układzie współrzędnych świata,
    \item wylosowania parametrów wewnętrznych kamery oraz parametrów dystorsji,
    \item wylosowania pozy i odległości szachownicy względem kamery,
    \item rzutowania punktów wzorca na płaszczyznę obrazu oraz zdeformowania obrazu zgodnie
          z modelem dystorsji optycznej.
\end{enumerate}

W przypadku datasetu sekwencyjnego jeden zestaw parametrów kamery jest stały dla całej
sekwencji, natomiast dla kolejnych klatek zmienia się położenie płaszczyzny szachownicy.
Taka konstrukcja dobrze odpowiada rzeczywistej procedurze kalibracji: kamera pozostaje tą
samą kamerą, a obserwowany wzorzec jest pokazywany pod różnymi kątami i w różnych miejscach
kadru. Model uczący się na takim zbiorze nie powinien zapamiętywać jednej konkretnej
perspektywy, lecz powinien rozpoznawać te cechy obrazu, które wynikają ze stałych parametrów
optycznych.

\section{Reprezentacja idealnej szachownicy}
Punktem wyjścia jest idealny obraz wzorca kalibracyjnego, czyli szachownicy złożonej
z naprzemiennych czarnych i białych pól. Matematycznie wzorzec ten jest traktowany jako
płaska figura leżąca na płaszczyźnie $Z=0$ w pewnym lokalnym układzie współrzędnych. Jeżeli
szachownica ma $n_x$ punktów charakterystycznych w poziomie oraz $n_y$ punktów charakterystycznych
w pionie, a rozmiar pojedynczego pola wynosi $s$, to punkty siatki można zapisać jako
\begin{equation}
    \mathbf{P}_{ij} =
    \begin{bmatrix}
        X_i \\
        Y_j \\
        0
    \end{bmatrix},
    \qquad i = 0,\dots,n_x-1, \quad j = 0,\dots,n_y-1.
\end{equation}
Współrzędne $X_i$ oraz $Y_j$ są oddalone od siebie o stały krok $s$:
\begin{equation}
    X_i = \left(i - \frac{n_x-1}{2}\right)s,
    \qquad
    Y_j = \left(j - \frac{n_y-1}{2}\right)s.
\end{equation}
Odjęcie środka siatki nie zmienia geometrii samego wzorca, ale sprawia, że środek szachownicy
leży blisko początku układu współrzędnych. Jest to wygodne, ponieważ obroty wokół osi układu
lokalnego odpowiadają wtedy intuicyjnemu obracaniu całej planszy wokół jej środka, a nie wokół
jednego z narożników.

Rozmiar pola $s$ nie musi być wyrażony w centymetrach ani milimetrach. W datasetcie syntetycznym
jest on dobierany względem rozmiaru obrazu, np. jako ułamek większego wymiaru kadru. Jeżeli
szerokość obrazu wynosi $W$, a wysokość $H$, to naturalną skalą jest
\begin{equation}
    M = \max(W,H),
\end{equation}
a rozmiar pola może być opisany jako
\begin{equation}
    s = \alpha_s M,
    \qquad \alpha_s \in [0.08, 0.16].
\end{equation}
Taki zapis powoduje, że wielkość wzorca automatycznie dopasowuje się do rozdzielczości obrazu.
Nie chodzi tutaj o odtworzenie fizycznego rozmiaru tablicy, lecz o wygenerowanie różnych skal
obserwacji, które zmuszają model do uczenia się zjawisk geometrycznych, a nie jednej stałej
wielkości obiektu.

\section{Losowanie parametrów wewnętrznych kamery}
Syntetyczna kamera jest opisana standardowym modelem kamery otworkowej. Jej parametry
wewnętrzne tworzą macierz
\begin{equation}
    K =
    \begin{bmatrix}
        f_x & 0   & c_x \\
        0   & f_y & c_y \\
        0   & 0   & 1
    \end{bmatrix},
\end{equation}
gdzie $f_x$ i $f_y$ są ogniskowymi wyrażonymi w pikselach, a $(c_x,c_y)$ jest punktem
głównym obrazu. Wartości te są losowane z zadanych przedziałów, aby dataset zawierał wiele
różnych kamer, a nie tylko jedną konfigurację optyczną.

Ogniskowe są skalowane przez większy wymiar obrazu:
\begin{equation}
    f_x = \alpha_x M,
    \qquad
    f_y = \alpha_y M,
    \qquad
    \alpha_x,\alpha_y \in [0.8,1.4].
\end{equation}
Zakres ten oznacza, że kamera może mieć zarówno szersze, jak i węższe pole widzenia.
Mniejsza ogniskowa powoduje silniejszą perspektywę i większą podatność na widoczne
zniekształcenia geometryczne przy krawędziach kadru. Większa ogniskowa daje obraz bardziej
zbliżony do projekcji równoległej, w którym zmiany skali wynikające z głębi są mniej
gwałtowne.

Punkt główny jest losowany w pobliżu środka obrazu:
\begin{equation}
    c_x = \beta_x W,
    \qquad
    c_y = \beta_y H,
    \qquad
    \beta_x,\beta_y \in [0.45,0.55].
\end{equation}
Dopuszczenie przesunięcia punktu głównego jest ważne, ponieważ w rzeczywistych kamerach oś
optyczna nie musi przecinać matrycy dokładnie w geometrycznym środku. Nawet niewielka zmiana
$(c_x,c_y)$ wpływa na to, względem którego punktu liczona jest odległość radialna, a więc
bezpośrednio wpływa na przestrzenny rozkład dystorsji.

\section{Model dystorsji radialnej i tangencjalnej}
Poza idealną projekcją perspektywiczną obraz jest deformowany przez dystorsję optyczną.
W procesie generowania wykorzystany jest model Brown--Conrady, powszechnie stosowany
w kalibracji kamer i zgodny z konwencją używaną przez OpenCV. Model ten opisuje punkt
w znormalizowanych współrzędnych kamery jako
\begin{equation}
    x = \frac{u-c_x}{f_x},
    \qquad
    y = \frac{v-c_y}{f_y},
\end{equation}
gdzie $(u,v)$ oznacza współrzędne pikselowe punktu w obrazie, natomiast $(x,y)$ jest jego
odpowiednikiem w układzie znormalizowanym. Odległość punktu od osi optycznej wyraża się przez
\begin{equation}
    r^2 = x^2 + y^2.
\end{equation}

Część radialna dystorsji zależy wyłącznie od odległości od punktu głównego. Jej współczynnik
skali ma postać
\begin{equation}
    R(r) = 1 + k_1 r^2 + k_2 r^4 + k_3 r^6.
\end{equation}
Jeżeli dominujący współczynnik radialny powoduje zmniejszanie współrzędnych dla dużych $r$,
obraz przyjmuje charakter beczkowy: linie proste wyginają się na zewnątrz, szczególnie przy
brzegach kadru. Jeżeli efekt działa przeciwnie, powstaje dystorsja poduszkowa, w której linie
proste są wciągane w stronę środka. Z tego powodu losowanie wartości dodatnich i ujemnych jest
istotne: dataset obejmuje oba typy deformacji, a model nie uczy się tylko jednego znaku
zniekształcenia.

Oprócz dystorsji radialnej uwzględniona jest dystorsja tangencjalna. Wynika ona z nieidealnego
ustawienia soczewki względem płaszczyzny sensora i powoduje asymetryczne przesunięcia punktów.
Pełne równania zniekształcenia mają postać
\begin{align}
    x_d &= x\,R(r) + 2p_1xy + p_2(r^2 + 2x^2), \\
    y_d &= y\,R(r) + p_1(r^2 + 2y^2) + 2p_2xy.
\end{align}
Wektor współczynników dystorsji można więc zapisać jako
\begin{equation}
    \mathbf{d} = [k_1,k_2,p_1,p_2,k_3]^T.
\end{equation}
W syntetycznym datasetcie współczynniki są losowane z przedziałów
\begin{align}
    k_1 &\in [-0.7,0.7], &
    k_2 &\in [-0.4,0.4], &
    k_3 &\in [-0.2,0.2], \\
    p_1 &\in [-0.2,0.2], &
    p_2 &\in [-0.2,0.2].
\end{align}
Tak szerokie zakresy sprawiają, że część przykładów może zawierać bardzo wyraźne deformacje.
Jest to celowe z punktu widzenia uczenia nadzorowanego: silne przypadki dostarczają modelowi
wyraźnego sygnału, jak zmiana parametrów przekłada się na geometrię obrazu. Jednocześnie losowanie
ciągłe, a nie wybór kilku dyskretnych wartości, zmniejsza ryzyko, że sieć potraktuje problem
jak klasyfikację do kilku gotowych wariantów.

Po obliczeniu współrzędnych zniekształconych $(x_d,y_d)$ następuje powrót do pikselowego układu
obrazu:
\begin{equation}
    u_d = f_x x_d + c_x,
    \qquad
    v_d = f_y y_d + c_y.
\end{equation}
Końcowy obraz jest więc nie tylko obrazem płaskiej szachownicy po przekształceniu
perspektywicznym, lecz obrazem po przejściu przez nieliniowy model optyczny.

\section{Poza szachownicy względem kamery}
Każda klatka wymaga określenia położenia szachownicy w układzie kamery. Poza wzorca jest
opisana przez macierz obrotu $R$ oraz wektor translacji $\mathbf{t}$. Punkt wzorca
$\mathbf{P}$ w układzie lokalnym jest przekształcany do układu kamery zgodnie z równaniem
\begin{equation}
    \mathbf{P}_c = R\mathbf{P} + \mathbf{t}.
\end{equation}
Macierz $R$ powstaje ze złożenia trzech obrotów: pochylenia, odchylenia oraz obrotu w płaszczyźnie
obrazu. Jeżeli kąty oznaczymy jako $\theta$ dla pitch, $\psi$ dla yaw oraz $\phi$ dla roll,
to macierze elementarne można zapisać jako
\begin{equation}
    R_x(\theta) =
    \begin{bmatrix}
        1 & 0 & 0 \\
        0 & \cos\theta & -\sin\theta \\
        0 & \sin\theta & \cos\theta
    \end{bmatrix},
\end{equation}
\begin{equation}
    R_y(\psi) =
    \begin{bmatrix}
        \cos\psi & 0 & \sin\psi \\
        0 & 1 & 0 \\
        -\sin\psi & 0 & \cos\psi
    \end{bmatrix},
\end{equation}
\begin{equation}
    R_z(\phi) =
    \begin{bmatrix}
        \cos\phi & -\sin\phi & 0 \\
        \sin\phi & \cos\phi & 0 \\
        0 & 0 & 1
    \end{bmatrix}.
\end{equation}
Złożony obrót ma postać
\begin{equation}
    R = R_z(\phi)R_y(\psi)R_x(\theta).
\end{equation}

W datasetcie kąty są losowane w ograniczonych zakresach:
\begin{equation}
    \theta \in [-20^\circ,20^\circ],
    \qquad
    \psi \in [-20^\circ,20^\circ],
    \qquad
    \phi \in [-10^\circ,10^\circ].
\end{equation}
Takie przedziały pozwalają uzyskać widoczną zmianę perspektywy, ale nie prowadzą zwykle
do skrajnych przypadków, w których tablica jest niemal równoległa do kierunku patrzenia kamery
i traci czytelność. Z punktu widzenia kalibracji najcenniejsze są właśnie umiarkowanie różne
pozy: dostarczają informacji o geometrii, a jednocześnie zachowują wystarczającą ilość
widocznych narożników oraz struktur liniowych.

Translacja jest dana przez
\begin{equation}
    \mathbf{t} =
    \begin{bmatrix}
        t_x \\
        t_y \\
        t_z
    \end{bmatrix},
\end{equation}
gdzie składowe boczne są losowane względem szerokości i wysokości obrazu,
\begin{equation}
    t_x = \gamma_x W,
    \qquad
    t_y = \gamma_y H,
    \qquad
    \gamma_x,\gamma_y \in [-0.2,0.2],
\end{equation}
a odległość od kamery względem większego wymiaru kadru:
\begin{equation}
    t_z = \gamma_z M,
    \qquad
    \gamma_z \in [0.9,1.8].
\end{equation}
Parametr $t_z$ kontroluje skalę widocznej szachownicy. Mała odległość sprawia, że wzorzec
zajmuje większą część kadru, a perspektywa jest bardziej wyraźna. Większa odległość zmniejsza
wzorzec i osłabia różnice skali między bliższą i dalszą częścią płaszczyzny.

\section{Projekcja perspektywiczna}
Po przekształceniu punktu szachownicy do układu kamery otrzymujemy
\begin{equation}
    \mathbf{P}_c =
    \begin{bmatrix}
        X_c \\
        Y_c \\
        Z_c
    \end{bmatrix}.
\end{equation}
Idealna projekcja perspektywiczna, jeszcze bez dystorsji, sprowadza punkt trójwymiarowy
do znormalizowanych współrzędnych obrazu:
\begin{equation}
    x = \frac{X_c}{Z_c},
    \qquad
    y = \frac{Y_c}{Z_c}.
\end{equation}
Następnie macierz $K$ przekształca je do współrzędnych pikselowych:
\begin{equation}
    \begin{bmatrix}
        u \\
        v \\
        1
    \end{bmatrix}
    \sim
    K
    \begin{bmatrix}
        x \\
        y \\
        1
    \end{bmatrix}.
\end{equation}
W zapisie bezpośrednim daje to równania
\begin{equation}
    u = f_x\frac{X_c}{Z_c} + c_x,
    \qquad
    v = f_y\frac{Y_c}{Z_c} + c_y.
\end{equation}

Ponieważ szachownica jest płaska, całe przekształcenie pomiędzy jej lokalną płaszczyzną
a płaszczyzną obrazu można opisać homografią. Dla punktu płaszczyzny zapisanego jako
$(X,Y,1)^T$ zachodzi
\begin{equation}
    \lambda
    \begin{bmatrix}
        u \\
        v \\
        1
    \end{bmatrix}
    =
    H
    \begin{bmatrix}
        X \\
        Y \\
        1
    \end{bmatrix},
\end{equation}
gdzie $\lambda$ jest niezerowym współczynnikiem skali, a homografia ma postać
\begin{equation}
    H = K\,[\mathbf{r}_1\;\mathbf{r}_2\;\mathbf{t}],
\end{equation}
przy czym $\mathbf{r}_1$ oraz $\mathbf{r}_2$ są dwiema pierwszymi kolumnami macierzy obrotu
$R$. Trzecia kolumna $R$ nie pojawia się wprost, ponieważ wszystkie punkty wzorca mają
$Z=0$. Jest to klasyczny wynik geometrii projekcyjnej: obraz dowolnej płaszczyzny w kamerze
perspektywicznej jest związany z płaszczyzną źródłową homografią.

Homografia odpowiada za to, że równoległe linie szachownicy w obrazie źródłowym mogą zbiegać
się w punkcie zbiegu, a pola położone dalej od kamery mogą być widocznie mniejsze niż pola
bliższe. Bez tego etapu dataset zawierałby jedynie płaskie, frontalne obrazy wzorca, które
nie byłyby wystarczająco informacyjne dla kalibracji.

\section{Złożenie homografii i dystorsji}
Końcowy obraz powstaje jako złożenie projekcji płaszczyzny oraz nieliniowego odwzorowania
dystorsyjnego. Można to zapisać symbolicznie jako
\begin{equation}
    I_{\text{out}} = \mathcal{D}_{K,\mathbf{d}}\left(\mathcal{H}_{K,R,\mathbf{t}}(I_{\text{board}})\right),
\end{equation}
gdzie $I_{\text{board}}$ oznacza idealny obraz szachownicy, $\mathcal{H}_{K,R,\mathbf{t}}$
jest przekształceniem perspektywicznym wynikającym z macierzy kamery i pozy tablicy,
a $\mathcal{D}_{K,\mathbf{d}}$ jest deformacją optyczną zadaną przez współczynniki
$\mathbf{d}$.

Ważne jest rozróżnienie tych dwóch zjawisk. Homografia opisuje zmianę wyglądu płaskiego
obiektu wynikającą z jego ustawienia względem kamery. Gdyby kamera była idealna, proste linie
szachownicy pozostałyby prostymi liniami po projekcji perspektywicznej. Dystorsja natomiast
powoduje, że nawet obraz linii prostych może stać się krzywoliniowy. Dlatego dla uczenia
parametrów kamery najważniejsze są subtelne różnice pomiędzy deformacją perspektywiczną
a deformacją optyczną. Perspektywa zależy od pozy tablicy w danej klatce, natomiast dystorsja
jest własnością kamery i pozostaje taka sama dla wszystkich obserwacji w sekwencji.

Z praktycznego punktu widzenia obraz wynikowy musi zostać obliczony na dyskretnej siatce pikseli.
Wartość piksela po przekształceniu nie zawsze odpowiada dokładnie jednemu pikselowi obrazu
źródłowego, dlatego stosowana jest interpolacja. Matematycznie można traktować obraz jako
funkcję
\begin{equation}
    I: \Omega \subset \mathbb{R}^2 \rightarrow \mathbb{R}^3,
\end{equation}
która dla dowolnych współrzędnych ciągłych zwraca przybliżoną wartość koloru. Interpolacja
liniowa jest kompromisem: zachowuje płynność przekształcenia i nie wprowadza tak ostrych
artefaktów jak najbliższy sąsiad, a jednocześnie jest prostsza od interpolacji wyższych rzędów.

\section{Kontrola widoczności wzorca}
Losowanie parametrów mogłoby prowadzić do obrazów, w których szachownica prawie całkowicie
wychodzi poza kadr albo jest widoczna jedynie jako mały fragment. Takie przykłady są zwykle
mało przydatne, ponieważ nie zawierają wystarczającej informacji geometrycznej. Dlatego proces
generowania ogranicza przyjmowane transformacje: po wylosowaniu pozy sprawdzane jest, czy
rzutowane punkty lub istotne narożniki mieszczą się w granicach obrazu.

Warunek widoczności można zapisać jako
\begin{equation}
    0 \leq u_i < W,
    \qquad
    0 \leq v_i < H
\end{equation}
dla wybranych punktów kontrolnych $i$. Jeżeli warunek nie jest spełniony, próbka może zostać
odrzucona i zastąpiona nowym losowaniem. Taka procedura nie zmienia matematycznego modelu
kamery, ale zmienia rozkład danych: usuwa przypadki skrajnie nieczytelne i koncentruje zbiór
na obserwacjach, które rzeczywiście mogłyby posłużyć do kalibracji.

Jednocześnie zbyt agresywne ograniczenie widoczności mogłoby prowadzić do zbyt łatwego datasetu.
Jeżeli wszystkie szachownice byłyby idealnie wyśrodkowane i niemal frontalne, model mógłby
osiągać dobre wyniki bez rozumienia pełnej geometrii problemu. Dlatego kontrola widoczności
powinna usuwać przypadki bezużyteczne, ale pozostawiać przesunięcia, obroty i zmiany skali,
które są potrzebne do nauczenia odporności.

\section{Budowa sekwencji obrazów}
Wariant sekwencyjny datasetu ma szczególne znaczenie, ponieważ naśladuje rzeczywistą sytuację,
w której jedna kamera obserwuje tę samą tablicę kalibracyjną w wielu położeniach. Dla pojedynczej
sekwencji losowany jest stały zestaw parametrów
\begin{equation}
    \Theta = (f_x,f_y,c_x,c_y,k_1,k_2,p_1,p_2,k_3,s,n_x,n_y),
\end{equation}
natomiast dla każdej klatki $t$ losowana jest osobna poza
\begin{equation}
    \Pi_t = (R_t,\mathbf{t}_t).
\end{equation}
Klatka sekwencji powstaje więc jako
\begin{equation}
    I_t = \mathcal{D}_{K,\mathbf{d}}\left(\mathcal{H}_{K,R_t,\mathbf{t}_t}(I_{\text{board}})\right),
    \qquad t=1,\dots,T.
\end{equation}
Etykieta pozostaje wspólna dla całej sekwencji:
\begin{equation}
    \mathbf{y}_{\text{seq}} = [f_x,f_y,c_x,c_y,k_1,k_2,p_1,p_2,k_3]^T.
\end{equation}

Ta konstrukcja wymusza istotną własność uczenia. Jeżeli model analizuje tylko pojedynczą
klatkę, musi rozdzielić wpływ dwóch czynników: aktualnego ustawienia szachownicy i stałych
parametrów kamery. W sekwencji wiele różnych póz ma te same parametry wewnętrzne i tę samą
dystorsję. Informacja powtarzalna między klatkami jest więc kandydatem na własność kamery,
a informacja zmienna między klatkami jest związana głównie z położeniem tablicy. Z tego powodu
sekwencja jest bogatszym sygnałem treningowym niż pojedynczy obraz.

Można to ująć probabilistycznie. Dla jednej sekwencji parametry kamery są losowane raz:
\begin{equation}
    \Theta \sim p(\Theta),
\end{equation}
a pozy kolejnych klatek są losowane warunkowo niezależnie:
\begin{equation}
    \Pi_t \sim p(\Pi),
    \qquad t=1,\dots,T.
\end{equation}
Obrazy są następnie deterministycznym wynikiem funkcji generującej:
\begin{equation}
    I_t = G(\Theta,\Pi_t).
\end{equation}
Zadanie modelu uczącego można zapisać jako znalezienie funkcji $F$, która dla sekwencji obrazów
odtwarza parametry kamery:
\begin{equation}
    F(I_1,\dots,I_T) \approx [f_x,f_y,c_x,c_y,k_1,k_2,p_1,p_2,k_3]^T.
\end{equation}

\section{Znaczenie etykiet i normalizacji}
Wygenerowany obraz sam w sobie nie wystarcza do uczenia nadzorowanego. Każda próbka musi mieć
przypisaną etykietę, czyli wektor wartości, które model ma przewidywać. W przypadku kalibracji
etykietą nie jest klasa ani opis tekstowy, lecz wektor ciągłych parametrów geometrycznych.
Najważniejszy wektor etykiety ma postać
\begin{equation}
    \mathbf{y} = [f_x,f_y,c_x,c_y,k_1,k_2,p_1,p_2,k_3]^T.
\end{equation}

Parametry w tym wektorze mają różne skale. Ogniskowe i punkt główny są mierzone w pikselach
i mogą mieć wartości rzędu setek, natomiast współczynniki dystorsji są bezwymiarowe i zwykle
mają wartości znacznie bliższe zeru. Gdyby trenować model bez uwzględnienia tej różnicy, błąd
ogniskowej mógłby numerycznie dominować nad błędem dystorsji. Z tego powodu naturalnym krokiem
jest normalizacja parametrów pikselowych przez rozmiar obrazu:
\begin{equation}
    \tilde{f}_x = \frac{f_x}{W},
    \qquad
    \tilde{f}_y = \frac{f_y}{H},
    \qquad
    \tilde{c}_x = \frac{c_x}{W},
    \qquad
    \tilde{c}_y = \frac{c_y}{H}.
\end{equation}
Można również normalizować ogniskowe przez $M=\max(W,H)$, jeżeli taki sam czynnik skali był
używany podczas losowania. Najważniejsza zasada jest taka, aby etykiety wejściowe do funkcji
straty miały porównywalne zakresy i aby odwrotne przeskalowanie było jednoznaczne podczas
interpretacji wyniku.

\section{Dlaczego dataset syntetyczny jest użyteczny dla kalibracji}
Największą zaletą syntetycznego datasetu jest pełna kontrola nad parametrami generującymi.
W rzeczywistym zbiorze danych dokładne wartości $f_x$, $f_y$, $c_x$, $c_y$ i współczynników
dystorsji trzeba najpierw oszacować tradycyjną metodą kalibracyjną, co samo wprowadza błąd
pomiarowy. W zbiorze syntetycznym etykieta jest znana, ponieważ to ona była przyczyną
wygenerowania obrazu. Pozwala to tworzyć bardzo duże zbiory danych bez ręcznej kalibracji
każdej konfiguracji.

Drugą zaletą jest możliwość świadomego pokrycia przestrzeni parametrów. Można wygenerować
kamery o różnych ogniskowych, różnych przesunięciach punktu głównego oraz dodatniej i ujemnej
dystorsji radialnej. W danych rzeczywistych wiele takich przypadków byłoby trudnych do zebrania,
bo wymagałoby wielu fizycznych kamer lub wielu obiektywów. Tutaj zmiana kamery sprowadza się
do zmiany wektora parametrów, a więc do kontrolowanego losowania z ustalonego rozkładu.

Trzecią zaletą jest możliwość tworzenia sekwencji, w których znana jest zależność między
klatkami. W każdej sekwencji zmienna jest perspektywa, ale stałe są parametry optyczne.
Jest to bardzo ważne dla modeli sekwencyjnych, ponieważ dostają one kilka obserwacji tego
samego układu optycznego. W takim ustawieniu model może uczyć się rozpoznawać parametry
kamery jako cechę wspólną, a nie przypadkowy efekt jednej pozy wzorca.

\section{Ograniczenia procesu generowania}
Mimo wielu zalet dataset syntetyczny nie jest idealnym odpowiednikiem danych rzeczywistych.
Opisany proces koncentruje się przede wszystkim na geometrii: projekcji perspektywicznej,
pozie płaszczyzny, ogniskowej, punkcie głównym i dystorsji. Rzeczywiste obrazy zawierają jednak
również szum sensora, nierównomierne oświetlenie, rozmycie ruchu, winietowanie, kompresję,
odbicia światła, niedoskonałości wydruku oraz lokalne deformacje fizycznej tablicy.
Jeżeli tych efektów nie ma w danych syntetycznych, model może osiągać bardzo dobre wyniki
na wygenerowanych obrazach, ale gorzej przenosić się na obrazy z prawdziwej kamery.

Ograniczeniem jest również przyjęty zakres losowania. Model uczony tylko na kamerach z
ogniskowymi z przedziału $[0.8M,1.4M]$ może słabiej radzić sobie z kamerą spoza tego zakresu.
Podobnie, jeżeli pozycje tablicy są ograniczone do umiarkowanych kątów, model nie musi poprawnie
obsługiwać bardzo ostrych perspektyw. Dlatego zakresy generowania należy traktować jako część
definicji problemu: określają one, jaką rodzinę kamer i obserwacji ma obejmować wytrenowany
model.

\section{Podsumowanie procesu}
Proces tworzenia datasetu można podsumować jako kontrolowane losowanie parametrów kamery
i pozy wzorca, a następnie wygenerowanie obrazu przez model geometryczny. Dla każdej próbki
powstaje najpierw idealna płaska szachownica, następnie wybierane są parametry macierzy
wewnętrznej $K$, współczynniki dystorsji $\mathbf{d}$ oraz poza $(R,\mathbf{t})$.
Szachownica jest rzutowana na obraz przez homografię wynikającą z modelu kamery otworkowej,
a potem deformowana przez równania Brown--Conrady'ego.

W wariancie sekwencyjnym ten sam zestaw parametrów kamery obowiązuje dla wielu klatek,
natomiast każda klatka otrzymuje inną pozę wzorca. Dzięki temu jedna sekwencja przedstawia
wiele obserwacji tej samej kamery. Etykietą sekwencji jest wektor parametrów
\begin{equation}
    [f_x,f_y,c_x,c_y,k_1,k_2,p_1,p_2,k_3]^T,
\end{equation}
czyli dokładnie te wartości, które model ma przewidywać. Matematycznie dataset jest więc
zbiorem par
\begin{equation}
    \left((I_1,\dots,I_T),\mathbf{y}_{\text{seq}}\right),
\end{equation}
gdzie obrazy są wynikiem znanego procesu projekcji i dystorsji, a etykieta jest parametrem
tego procesu. Taka konstrukcja sprawia, że zadanie uczenia jest dobrze określone: sieć ma
odwrócić, w przybliżeniu, syntetyczny proces generowania i z obrazu lub sekwencji obrazów
odzyskać parametry kamery.

\end{document}